\chapter{Estado del arte}
\section{Titulo nivel 2}
La globalización ha generado una transformación significativa en la forma en que las sociedades interactúan, comparten información y desarrollan sus economías. Esta interconexión ha traído consigo múltiples beneficios, como el acceso a nuevas tecnologías y mercados, pero también ha planteado desafíos importantes en términos de identidad cultural, sostenibilidad y equidad. Comprender estos efectos resulta clave para formular estrategias que permitan un desarrollo más equilibrado y justo.

\subsection{Ejemplo de ubicación de viñetas}
\begin{itemize}
    \item •	Comprender estos efectos resulta clave para formular estrategias que permitan un desarrollo más equilibrado y justo.
    \item •	Comprender estos efectos resulta clave para formular estrategias que permitan un desarrollo más equilibrado y justo.
    \item •	Comprender estos efectos resulta clave para formular estrategias que permitan un desarrollo más equilibrado y justo.
\end{itemize}
\subsection{Titulo nivel 3}
El uso de herramientas digitales en entornos educativos ha demostrado un potencial considerable para mejorar el aprendizaje, especialmente en contextos donde el acceso a recursos tradicionales es limitado. Sin embargo, la brecha digital persiste como una barrera que impide una participación equitativa. Es necesario que las políticas públicas y las instituciones académicas trabajen conjuntamente para garantizar una inclusión real, promoviendo tanto la infraestructura como la capacitación adecuada.
\subsubsection{Titulo nivel 4}
El uso de herramientas digitales en entornos educativos ha demostrado un potencial considerable para mejorar el aprendizaje, especialmente en contextos donde el acceso a recursos tradicionales es limitado. Sin embargo, la brecha digital persiste como una barrera que impide una participación equitativa. Es necesario que las políticas públicas y las instituciones académicas trabajen conjuntamente para garantizar una inclusión real, promoviendo tanto la infraestructura como la capacitación adecuada.

\paragraph{Titulo nivel 5}
Uno de los aspectos más relevantes al analizar cualquier fenómeno social es la necesidad de integrar diferentes perspectivas disciplinarias. La complejidad de los problemas actuales requiere enfoques interdisciplinarios que consideren factores económicos, culturales, ambientales y tecnológicos. Solo así es posible alcanzar una comprensión más profunda de la realidad y proponer soluciones sostenibles y efectivas.
La complejidad de los problemas actuales requiere enfoques interdisciplinarios que consideren factores económicos, culturales, ambientales y tecnológicos. Solo así es posible alcanzar una comprensión más profunda de la realidad y proponer soluciones sostenibles y efectivas.
Uno de los aspectos más relevantes al analizar cualquier fenómeno social es la necesidad de integrar diferentes perspectivas disciplinarias. La complejidad de los problemas actuales requiere enfoques interdisciplinarios que consideren factores económicos, culturales, ambientales y tecnológicos. Solo así es posible alcanzar una comprensión más profunda de la realidad y proponer soluciones sostenibles y efectivas.

\chapter{Figuras}

\section*{Ejemplo de Figura con su Nota}

\begin{figure}[htbp]
    \centering
    \includegraphics[width=0.75\textwidth]{image1.png}
    \caption{\textit{Pie de Figura}}
    \label{fig:ejemplo-nota}
\end{figure}

\textbf{Nota.} Se hace una descripción del contenido de la tabla en cuestión de lo que se esté exponiendo dentro de esta. Cuando la figura es de elaboración propia no es necesario agregar ningún tipo de declaración de derechos de autor. En APA se asume que todo lo que no tenga cita (o la declaración de derechos de autor) es de autoría del propio autor.

\chapter{Tablas}

\section*{}

\begin{table}[htbp]
     \caption{\textit{El Título Debe ser Breve y Descriptivo (en Cursiva)}}
    \centering
    \begin{tabular}{p{7cm} p{7cm}}
        \toprule
        \textbf{Columna uno} & \textbf{Columna dos} \\
        \midrule
        Table data & Table data \\
        Table data & Table data \\
        Table data & Table data \\
        Table data & Table data \\
        Table data & Table data \\
        \bottomrule
    \end{tabular}
\end{table}

% Nota APA debajo de la tabla, alineada a la izquierda
\noindent\textit{Nota.} Se hace una descripción del contenido de la tabla en cuestión de lo que se esté exponiendo dentro de esta, para referenciar la tabla se puede tomar el ejemplo de la figura que está ubicada en este documento.

